% BHU PYQ -- Auto-generated & error-corrected
\documentclass[11pt,a4paper]{article}

\usepackage[a4paper,top=2.5cm,bottom=2.5cm,left=2.8cm,right=2.8cm,headheight=14pt]{geometry}
\usepackage{amsmath,amssymb}
\usepackage[T1]{fontenc}
\usepackage[utf8]{inputenc}
\usepackage{lmodern}
\usepackage{microtype}
\usepackage{fancyhdr}
\usepackage{enumitem}
\usepackage{hyperref}
\usepackage{lastpage}
\usepackage{parskip}

\hypersetup{colorlinks=true,urlcolor=black,linkcolor=black,
  pdftitle={BSC-13A: Statistical Inference, Sampling and Design of Experiments -- BHU PYQ},pdfauthor={Banaras Hindu University}}

\pagestyle{fancy}\fancyhf{}
\renewcommand{\headrulewidth}{0.4pt}\renewcommand{\footrulewidth}{0.4pt}
\fancyhead[L]{\small   \textit{Banaras Hindu University}}
\fancyhead[R]{\small   \textit{BSC-13A: Inference, Sampling \& Design of Exp}}
\fancyfoot[C]{\small Page  \thepage\ of \pageref{LastPage}}


\newlist{parts}{enumerate}{2}
\setlist[parts,1]{label=(\alph*),leftmargin=*,itemsep=5pt,topsep=3pt}
\setlist[parts,2]{label=(\roman*),leftmargin=*,itemsep=3pt,topsep=2pt}

\newcommand{\pts}[1]{\hfill   \textit{[#1]}}


\begin{document}


\begin{center}
  {\large\bfseries BANARAS HINDU UNIVERSITY}\\[2pt]
  {
\normalsize B.Sc. (Hons.) Semester IV Examination, 2013-14}\\[6pt]
  \rule{0.8\linewidth}{0.5pt}\\[6pt]
  {\large\bfseries Paper No. BSC-13A: Statistical Inference, Sampling and Design of Experiments}\\[4pt]
  {
\normalsize Time: 3 Hours\hfill Maximum Marks: 70}
\end{center}
\rule{\linewidth}{0.4pt}
\smallskip



\noindent   \textit{Instructions: Answer   \textbf{five} questions including Question~1 which is compulsory. Symbols have their usual meanings.}
\bigskip

\medskip


\noindent   \textbf{Question 1.}

\begin{parts}
  \item Define probability. State and prove the addition theorem of probability for any two events $A$ and $B$, i.e.,
  \[
  P(A \cup B) = P(A) + P(B) - P(A \cap B)
  \]
  Also state the theorem if $A$ and $B$ are independent events. \pts{7}
  \item A balanced die is thrown two times. Find the probability that the sum of the numbers in the two throws is:
  
\begin{parts}
    \item 7
    \item more than 10
    \item less than 5
  \end{parts} \pts{7}
\end{parts}
\medskip\rule{\linewidth}{0.2pt}

\medskip


\noindent   \textbf{Question 2.}

\begin{parts}
  \item Define a discrete random variable and its probability mass function (p.m.f.) giving a suitable example. \pts{7}
  \item Explain the binomial and Poisson distributions. Discuss their properties and mention fields of their applications. \pts{7}
\end{parts}
\medskip\rule{\linewidth}{0.2pt}

\medskip


\noindent   \textbf{Question 3.}

\begin{parts}
  \item Explain the probability density function (p.d.f.) and cumulative distribution function (c.d.f.) of a continuous random variable. \pts{7}
  \item Check whether the function defined as:
  \[
  f(x) = 
\begin{cases}
    e^{-x}, & x > 0 \\
    0, &  \text{otherwise}
  \end{cases}
  \]
  is a probability density function. \pts{7}
\end{parts}
\medskip\rule{\linewidth}{0.2pt}

\medskip


\noindent   \textbf{Question 4.}

\begin{parts}
  \item State the important properties of a normal distribution. \pts{7}
  \item Explain the concept of testing of hypothesis. In this context, explain two types of errors (Type I and Type II errors) and the level of significance. \pts{7}
\end{parts}
\medskip\rule{\linewidth}{0.2pt}

\medskip


\noindent   \textbf{Question 5.}

\begin{parts}
  \item Explain the analysis of variance (ANOVA) table for a one-way classified data (under Completely Randomized Design). \pts{7}
  \item Explain two uses of the $\chi^2$-test and two uses of the $t$-test. \pts{7}
\end{parts}
\medskip\rule{\linewidth}{0.2pt}

\medskip


\noindent   \textbf{Question 6.}

\begin{parts}
  \item Explain the terms: population, sample, and sampling frame giving suitable examples. \pts{7}
  \item Describe simple random sampling. Explain a method to select a simple random sample from a population. \pts{7}
\end{parts}
\medskip\rule{\linewidth}{0.2pt}

\medskip


\noindent   \textbf{Question 7.}
Describe stratified random sampling. Derive the variance of the stratified sample mean estimator and explain proportional and optimum (Neyman) allocations. \pts{14}
\medskip\rule{\linewidth}{0.2pt}

\medskip


\noindent   \textbf{Question 8.}
Write short notes on any   \textbf{three} of the following:

\begin{parts}
  \item Conditional probability
  \item Statistical definition of probability
  \item Mean of binomial and Poisson distributions
  \item Large sample test for single mean with one example
  \item Advantages of sampling over complete enumeration
\end{parts}
\pts{14}

\vfill
\rule{\linewidth}{0.4pt}

\begin{center}\small
  \href{https://bhu-exam.vercel.app/}{Click here to visit our website}\\[4pt]
  \href{https://me-aryan.vercel.app/}{Click here to visit our contributor}\\[4pt]
  \href{https://quashan-me.vercel.app/}{Click here to visit our contributor}
\end{center}

\end{document}